% ===========================================================================
\section{Emission Factors, and a Trap That Leaves \COtwo{} Correct}
\label{sec:factors}
% ===========================================================================

\subsection{Which table, and why it matters here}

IPCC Volume~2 Chapter~3 publishes separate default tables for road, rail and off-road mobile
sources. Across them the \COtwo{} factor for a given fuel is identical, because \COtwo{}
depends only on the carbon in the fuel and on complete oxidation. The \CHfour{} and \NtwoO{}
factors are not.

\begin{table}[htbp]
\centering
\footnotesize
\begin{tabular}{@{}llrr@{}}
\toprule
\textbf{Category} & \textbf{Table} & \textbf{\COtwo{}} & \textbf{\NtwoO{}} \\
 & & \multicolumn{2}{c}{kg/TJ} \\
\midrule
Road diesel        & 3.2.2 & 74\,100 & 3.9 \\
Off-road machinery & 3.3.1 & 74\,100 & 28.6 \\
Railways           & 3.4.1 & 74\,100 & 28.6 \\
\bottomrule
\end{tabular}
\caption{The same diesel in three tables. \COtwo{} is identical by construction; \NtwoO{}
differs by a factor of seven, because off-road and locomotive engines generally carry no
catalyst and run lean at high load, which is the \NtwoO{}-forming regime.}
\label{tab:tables}
\end{table}

\subsection{What the wrong table costs, quantified}

Off-road machinery must take Table~3.3.1. If it took the road table instead, the engine
reports the diagnostic on every run: off-road \NtwoO{} would be \SI{0.022}{\mega\tonne}
instead of \SI{0.161}{\mega\tonne}, and the category total \SI{1.59}{\mega\tonne} instead of
\SI{1.73}{\mega\tonne}---\SI{8}{\percent} low.

The reason this is worth a section rather than a footnote is the failure signature. \COtwo{}
is \SI{91}{\percent} of the off-road category and is \emph{identical} between the two tables,
so an inventory making this error would still reconcile against the national diesel supply,
still agree with EDGAR to within its usual margin, and still pass every mass-balance check the
compiler had. Nothing available as a validation would catch it. The only thing that catches it
is reading the guidelines carefully enough to notice that there is more than one table---
precisely the step that time pressure removes.

It is also directional. Both the road and energy-industry tables carry the \emph{lower} trace-gas
values and are the more familiar ones, so a community making this error at random rates would
be biased low on non-\COtwo{} transport emissions rather than scattered around the truth.

A regression test pins the factor, so substituting the road value fails the build.

\subsection{Fossil methane, and why it is flagged}

Road and off-road diesel and petrol are fossil fuels, so their \CHfour{} takes the AR5 fossil
GWP of \num{30} rather than the biogenic \num{28}. Every methane-emitting component in this
inventory declares which it is, and a contract test asserts both that the declaration is
present and that the \COe{} conversion used the corresponding value.

The difference is small here---methane is \SI{0.40}{\mega\tonne} \COe{} across the
sector---but the flag is not optional, because the same field distinguishes dung-cake methane
from natural-gas methane in the residential sector where it matters considerably more.

\subsection{No separate oxidation factor}

The 2006 Guidelines assume complete oxidation and fold it into the emission factor;
Table~1.4 carries an explicit oxidation column of \num{1.0} for every fuel. Applying the
1996-vintage \num{0.98} on top would double-discount. This inventory's waste sector
\emph{does} carry an oxidation factor, because open burning genuinely is incomplete, and a
test pins both conventions so neither can be "corrected" into the other.
